\section{引言}

万历七年，张居正下令关闭天下书院六十四处；十四年后，顾宪成被罢，东林书院的讲学又进入政治斗争叙事。若只把前者叫作“收编”、后者叫作“突围”，论文仍停在漂亮的标签上。真正可检验的问题是：思想人物在关系网络中依靠何种边扩散，又通过何种边进入朝廷冲突？思想事件在时间上是连续积累，还是在少数节点被叙事压缩？

本文把《明朝那些事儿》这一通俗史著抽取成的知识图谱作为研究对象。我们不从哲学命题真伪出发，而从叙事怎样组织人物、事件与关系出发。原文检索仅用于确认语汇确实出现；关系类别构成衡量思想传播与权力交涉的结构位置；带年份事件则给出时间序列。三类证据各自回答“说过什么”“连到谁”“何时发生”，避免用一个含混的“思想影响力”代替全部现象。

本文贡献有三点。第一，提出可复算的操作化方案：以“师生同门”作为传播替代量，以君臣从属、政治攻讦、敌对冲突、同盟友党四类边作为权力接合替代量，并把自动推导的同事件边单独剔除。第二，给出思想型人物与权力型人物的可比构成：10名思想相关人物的显式关系中，师生同门占30.0\%，权力关系占58.9\%；张居正与严嵩两名参照人物的相应比例为11.7\%与80.5\%。第三，把思想相关事件排成年表，识别1508至1528年的形成—传授密集段，以及1579和1593年的制度性冲突节点。

本文的结论有明确边界。“突围/收编”是解释框架，不是图谱字段；关系类别与年份序列只能作为替代量。语料是通俗叙事而非《传习录》、奏疏或书院档案，因而本文能够说明这部书怎样叙述心学与朝廷，不能据此复原明代思想的全部内涵，也不能把词频当作社会影响力。
